\section{Diffentially Private Online-to-Batch}
\label{sec:DP-OTB_OTB}

In this section, we present our main differentially private online-to-batch algorithm. As a brief recap, online convex optimization (OCO) is a ``game'' in which for $T$ rounds, $t=1,\ldots,T$, an algorithms predicts a parameter $\vecw_t\in \calX$. It then receives a convex loss $\ell_t:\calX\to \RR$ and pays the loss $\ell_t(\vecw_t)$. The quality of the algorithm is measured by the regre defined as $\regret_T(\vecu)=\sum_{t=1}^T \ell_t(\vecw_t)-\ell_t(\vecu)$.

Online-to-batch algorithms \citep{cesa2004generalization} convert OCO algorithms (online learners) into stochastic convex optimization algorithms. In particular, for $\weight_1,\ldots,\weight_T>0$, the anytime online-to-batch conversion \citep{cutkosky2019anytime} defines the $t$-th loss as $\ell_t(\vecw)=\langle \weight_t\nabla F(\vecx_t,\xi_t), \vecw\rangle$, where $\vecx_t=\sum_{i=1}^t \frac{\weight_i}{\weight_{1:t}}\cdot \vecw_i$. Convergence of anytime online-to-batch builds on the following key result, and its proof is in Appendix \ref{app:DP-OTB:proof-A} for completeness.

\begin{restatable}[\cite{cutkosky2019anytime}]{lemma}{LemAnytimeOTB}
\label{lem:DP-OTB:basic-lemma}
For any sequence of $\weight_t>0, \vecg_t\in\RR^d$, suppose an online learner predicts $\vecw_t$ and receives $t$-th loss $\ell_t(\vecw)=\langle \vecg_t, \vecw\rangle$. Define $\vecx_t=\sum_{i=1}^t \frac{\weight_i}{\weight_{1:t}}\cdot\vecw_i$. Then for any convex and differentiable $f$,
\begin{align*}
    \weight_{1:T}(f(\vecx_T)-f(\vecu)) \leq \regret_T(\vecu) + \sum_{t=1}^T \langle \weight_t\nabla f(\vecx_t) - \vecg_t, \vecw_t-\vecu\rangle, \ \forall \vecu\in \RR^d.
\end{align*}
\end{restatable}

A tighter bound is possible \citep{joulani2020simpler}, but the simpler expression above suffices for our analysis.
As an immediate result, choosing $\weight_t=1$ and $\vecg_t=\weight_t\nabla F(\vecx_t,\xi_t)$ satisfies $\E[\vecg_t|\vecx_t]=\weight_t\nabla f(\vecx_t)$, so $\E[f(\vecx_T)-f(\vecu)]\le \E[\regret_T(\vecu)]/T$. Therefore, any online learner with sub-linear regret guarantees convergence for the last iterate $\vecx_T$. Notice that due to the formulation of the anytime online-to-batch, the iterate $\vecx_t$ is \emph{stable}. Consider the case where $\weight_t = 1$ for all $t$. Then, we have $\|\vecx_t-\vecx_{t-1}\|= \|\vecw_t-\vecx_{t-1}\|/t\le D/t$ \emph{regardless of what the online learner does}. This guarantee is significantly stronger than the classical online-to-batch result of \cite{cesa2004generalization}. We would like to take advantage of this stability to design our private online-to-batch algorithm. Intuitively, the algorithm has much lower sensitivity due to the stability of the iterates, which can be exploited to improve privacy.

Our goal will be to make the composition of $\vecg_1,\dots,\vecg_T$ private, which, in turn, makes the entire algorithm private since the algorithm is a post-processing of all gradients. To do so, we would like to add noise to the $\vecg_t$'s while still having $\vecg_t$ be a good estimate of $\nabla f(x_t)$. 
In standard non-private online-to-batch, $\vecg_t$ is usually defined as the stochastic gradient $\nabla F(x_t,z_t)$. However, directly adding noise to $\nabla F(x_t,z_t)$ is not a good idea because its sensitivity is $O(1)$ (more specifically the sensitivity is bounded by $2G$ by Lipschitzness). Consequently, we need to add noise to $\vecg_t$ whose noise (standard deviation) is of order $O(1/\rho)$ in order to achieve $(\alpha, \alpha\rho^2/2)$-RDP guarantee.

Instead, we design $\vecg_t$ to be a streaming sum of gradient difference estimators, i.e., $\vecg_t=\sum_{i=1}^t \vecd_i$, where $\vecd_i=\weight_i\nabla F(\vecx_i,\xi_i)-\weight_{i-1}\nabla F(\vecx_{i-1},\xi_i)$ and $\weight_0\equiv 0$. Note that both gradients are evaluated at the same data batch $\xi_i$.
As a brief illustration why this design helps improve privacy, let's assume $F$ is smooth and let $\weight_i=1$. Then $\|\vecd_i\|_*\le H\|\vecx_i-\vecx_{i-1}\|\le DH/i$ (recall the stability property that $\|\vecx_i-\vecx_{i-1}\|\le D/i$). In other words, the sensitivity of the gradient difference estimator is $O(1/i)$. Together with a carefully designed tree mechanism, the total noise could be reduced from $O(1/\rho)$ to $\tilde O(1/\rho t)$, thus allowing to achieve the optimal rate.

On the other hand, the tree mechanism makes $\vecg_t$ no longer an conditionally unbiased estimator of $\weight_t\nabla f(\vecx_t)$. More specifically, although it remains the case that $\E_{\xi_i}[\vecd_i|\xi_1,\ldots,\xi_{i-1}]=\weight_i\nabla f(\vecx_i)-\weight_{i-1}\nabla f(\vecx_{i-1})$, it is not necessarily true that $\E[\vecg_t|\xi_1,\ldots,\xi_{t-1}]= \weight_t\nabla f(x_t)$. Unbiasedness plays a key role in standard convergence analyses, but we will need a much more delicate analysis. 

Moreover, although we mostly discuss the $\weight_t =1$ case above for intuition, our algorithm is analyzed using the general case $\weight_t = t^k$ for $k\ge 1$. The guiding principle for this formula is the sensitivity of $\vecd_t$. For $k\ge 1$, the sensitivity of $\vecd_t$ is of order $O(t^{k-1})$. For $k=1$, this is a constant sensitivity, which is particularly intuitive for analysis. For $k=0$ (i.e. the standard weighting in online-to-batch), the sensitivity is actually $O(1/t)$, which is much more complicated to analyze. In order to apply the tree aggregation easily, we want the sensitivity of $\vecd_t$ to be polynomial in $t$, rather than the inverse polynomial $1/t$, so we define $\weight_t=t^k$ and ask $k\ge 2$.
Furthermore, in all cases except for the parameter-free case (Section \ref{sec:DP-OTB_parameter-free}), our results hold for all $k\geq 1$. In the parameter-free case, we choose $k=3$ for algebraic reasons.

The pseudo-code is presented in Algorithm \ref{alg:DP-OTB:main}, which has linear time complexity. It is similar to anytime online-to-batch, while we replaced $\vecg_t$ with the more complicated streaming sum estimator with added noise from the tree mechanism.

\begin{algorithm}[t]
    \begin{algorithmic}[1]
        \Statex \textbf{Input:} OCO algorithm $\calA$ with domain $\calW$, weights $\weight_1,\ldots,\weight_T$, data $\xi_1,\ldots,\xi_T$.
        \State \textbf{Initialize:} Set $\weight_0 = 0, \vecx_0 = 0$ and $\vecg_0 = 0$. 
        \For{$t=1,\ldots,T$}
            \State Get $\vecw_t$ from $\calA$ and compute $\vecx_t = (\weight_{1:t-1}\vecx_{t-1}+\weight_t \vecw_t)/\weight_{1:t}$.
            \State Compute gradient difference $\vecd_t = \weight_t\nabla F(\vecx_t,\xi_t)-\weight_{t-1}\nabla F(\vecx_{t-1},\xi_t)$.
            \State Update $\vecg_t = \vecg_{t-1} + \vecd_t$.
            \State Apply tree mechanism $\tilde\vecg_t = \vecg_t + \overline\vecr_t$, where $\overline\vecr_t=\tree(t)$.
            \State Send loss $\ell_t(\vecw) = \langle \tilde\vecg_t, \vecw\rangle$ to $\calA$ as the $t$-th loss.
        \EndFor
    \end{algorithmic}
    \caption{Differentially-Private Online-to-Batch}
    \label{alg:DP-OTB:main}
\end{algorithm}

\paragraph{Convergence}

Following Theorem \ref{lem:DP-OTB:basic-lemma} and Fenchel-Young's inequality, Algorithm \ref{alg:DP-OTB:main} satisfies:
\begin{align}
    \label{eq:DP-OTB:basic-decomposition}
    \weight_{1:T}(f(\vecx_T)-f(\vecx_\ast)) 
    &\leq \regret_T(\vecx_\ast) + \sum_{t=1}^T \langle \weight_t\nabla f(\vecx_t) - \vecg_t - \overline\vecr_t, \vecw_t - \vecx_\ast\rangle \\
    &\leq \regret_T(\vecx_\ast) + \sum_{t=1}^T D\|\weight_t\nabla f(\vecx_t)-\vecg_t\|_\ast + \sum_{t=1}^T D\|\overline\vecr_t\|_\ast.
    \label{eq:DP-OTB:basic-decomposition-2}
\end{align}

\begin{lemma}[DP-OTB Variance Bound]
\label{lem:DP-OTB:variance-expectation-bound}
If Assumptions \ref{as:DP-OTB:sc-norm}--\ref{as:DP-OTB:sigma-H} hold and $\weight_t = t^k$, then
\begin{equation}
    \E\left[\|\weight_t\nabla f(\vecx_t)-\vecg_t\|_\ast^2\right] \le 4(k+1)^2(\sigma_G^2+D^2\sigma_H^2)t^{2k-1}/\lambda.
    \label{eq:DP-OTB:variance-norm}
\end{equation}
\end{lemma}

As a remark, \eqref{eq:DP-OTB:basic-decomposition-2} decomposes the suboptimality gap $f(\vecx_T)-f(\vecx_\ast)$ into three components. The first term is the \emph{regret} of the online learner, $\regret_T(\vecx_\ast)$, which depends on the OCO subroutine and could not be simplified further. The second term captures the total \emph{bias error} of the biased sum of gradient difference estimator $\vecg_t$. Using a general martingale difference concentration inequality (formally stated as Theorem \ref{thm:concentration:martingale-difference}), we can show hat each error term is bounded by $\|\weight_t\nabla f(\vecx_t)-\vecg_t\|_*\approx t^{k-1/2}$. Lemma \ref{lem:DP-OTB:variance-expectation-bound} provides a more formal technical bound. Lastly, the third term captures the \emph{privacy noise} from the tree mechanism, and the previous result (Theorem \ref{thm:tree-mechanism:noise}) implies that $\|\overline\vecr_i\|_\ast \approx \RDPdist \cdot \sigma_t \approx \RDPdist \cdot t^{k-1}/\rho$, where $\sigma_t$'s are tunable noise levels in the tree mechanism whose optimal values are specified in Theorem \ref{thm:DP-OTB:main-privacy}. Overall, if we assume $\weight_t=1$ (i.e., $k=0$) and use the Gaussian mechanism with $\RDPdist=\sqrt{d}$, then \eqref{eq:DP-OTB:basic-decomposition-2} implies
\begin{equation}
\E[f(\vecx_T) - f(\vecx_\ast)] \approx \frac{1}{T}\E[\regret_T(\vecx_\ast)] + \frac{1}{\sqrt{T}} + \frac{\sqrt{d}}{\rho T}.
\label{eq:DP-OTB:main:approx-bound}
\end{equation}
We will provide a more careful analysis in the following section. Notably, for any online learner with $O(\sqrt{T})$ regret, the DP-OTB framework automatically achieves the \emph{near-optimal} convergence rate $\tilde O(1/\sqrt{T} + \sqrt{d} / \rho T)$.


\paragraph{Privacy}

Upon using the {\tree} mechanism with carefully chosen noises, we can ensure that DP-OTB (Algorithm \ref{alg:DP-OTB:main}) is RDP. The proof is deferred to Appendix \ref{app:DP-OTB:tree-aggregation-RDP}.

\begin{restatable}[DP-OTB Privacy]{theorem}{ThmDpOtbPrivacy}
\label{thm:DP-OTB:main-privacy}
Suppose $\|\cdot\|^2$ is $\lambda$-strongly convex, $\calX$ is bounded by $D$, $F$ is $G$-Lipschitz and $H$-smooth, and $\calR$ is a $(\alpha,\RDPdist)$-RDP distribution. 
Then Algorithm \ref{alg:DP-OTB:main} is $(\alpha,\alpha\rho^2/2)$-RDP if we define $\weight_t=t^k$ and $\sigma_t^2$ in {\tree} as 
\begin{equation}
    \sigma_t = 2(k+1)(G+DH)\sqrt{\log_2(2T)}t^{k-1} / \rho, 
    \label{eq:DP-OTB:RDP-noise}
\end{equation}
\end{restatable}

\paragraph{Main Result.} Now we can combine all the previous results to prove the privacy and convergence guarantee of our algorithm.

\begin{theorem}
\label{thm:DP-OTB:general-norm}
Suppose Assumption \ref{as:DP-OTB:sc-norm}--\ref{as:DP-OTB:sigma-H} hold, and $\calR$ is a $(\alpha,\RDPdist)$-RDP distribution. With $\weight_t=t^k$ and $\sigma_t^2$ defined in \eqref{eq:DP-OTB:RDP-noise}, DP-OTB is $(\alpha,\alpha\rho^2/2)$-RDP and 
% $\E[f(\vecx_T) - f(\vecx_\ast)]$ is bounded by:
\begin{align*}
    \E[f(\vecx_T) - f(\vecx_\ast)] \lesssim 
    \frac{\E[\regret_T(\vecx_\ast)]}{T^{k+1}}
    + \frac{D}{\sqrt{\lambda}}\left(\frac{\sigma_G+D\sigma_H}{\sqrt{T}}+\frac{\RDPdist(G+DH)\log(T)}{\rho T}\right).
    % NOTE: exact bound
    % \frac{(k+1)\E[\regret_T(\vecx_\ast)]}{T^{k+1}}
    % + \frac{2(k+1)^2D}{\sqrt{\lambda}}\left(\frac{\sigma_G+D\sigma_H}{\sqrt{T}}+\frac{\sqrt{2\RDPdist}(G+DH)\log_2(2T)}{\rho T}\right).
\end{align*}
% Moreover, recall that the online learner receives $t$-th loss $\ell_t(w)=\langle g_t+\gamma_t,w\rangle$. It holds that
Moreover, the private gradient estimator $\tilde\vecg_t = \vecg_t + \tree(t)$ satisfies
\begin{align*}
    \E[\|\tilde\vecg_t\|_\ast^2] \lesssim t^{2k}\left(G^2 + \frac{1}{\lambda}\left(\frac{(\sigma_G+D\sigma_H)^2}{t}+\frac{\RDPdist^2(G+DH)^2\log(T)}{\rho^2t^2}\right) \right).
    % NOTE: exact bound
    % \E[\|g_t+\gamma_t\|_*^2] \leq 3t^{2k}\left(G^2 + \frac{4(k+1)^2}{\lambda}\left(\frac{(\sigma_G+D\sigma_H)^2}{t}+\frac{2V(G+DH)^2\log_2^2(2T)}{\rho^2t^2}\right) \right).
\end{align*}
\end{theorem}

Before proving this result, let's consider the setting where $\calX$ is the Euclidean norm and $\calR$ is the Gaussian distribution (which is an $(\alpha,\sqrt{d})$-RDP distribution). Many OCO algorithms, such as OSD, OMD, and FTRL, satisfy that if $\E[\|\nabla\ell_t(\vecw_t)\|_\ast^2]\le G_\ell^2$ for all $t$, then $\E[\regret_T(\vecx_\ast)] \leq O(DG_\ell\sqrt{T})$. Hence, Theorem \ref{thm:DP-OTB:general-norm} yields a convergence bound
\begin{equation}
    \E[f(x_T)-f(\vecx_\ast)] = O\left(\frac{D(G+D\sigma_H)}{\sqrt{T}} + \frac{\sqrt{d}D(G+DH)\log T}{\rho T}\right).
    \label{eq:DP-OTB:main:formal-bound}
\end{equation}
This bound is of $\tilde O(1/\sqrt{T}+\sqrt{d}/\rho T)$, which is equivalent to the standard $(\rho,\gamma)$-DP bound of order $\tilde O(1/\sqrt{T} + \sqrt{d\log(1/\gamma)}/\rho T)$. Notably, this matches the optimal rate for private stochastic optimization with convex and smooth losses (\cite{bassily2019private}).

\begin{proof}[Proof of Theorem \ref{thm:DP-OTB:general-norm}]
From Eq. \eqref{eq:DP-OTB:basic-decomposition-2}, we have:
\begin{align*}
    & \E[f(\vecx_T)-f(\vecx_\star)] 
    \leq \frac{1}{\weight_{1:T}} \E\left[ \regret_T(\vecx_\star) + D\sum_{t=1}^T (\|\weight_t\nabla f(\vecx_t)-\vecg_t\|_\ast + \|\overline\vecr_t\|_\ast) \right] \\
    \intertext{
    Upon applying the variance bound in \eqref{eq:DP-OTB:variance-norm} on the second term and Theorem \ref{thm:tree-mechanism:noise} on the third term, together with Jensen's inequality $\E[\|X\|_*]\leq \sqrt{\E[\|X\|_*^2]}$, we have}
    &\lesssim \frac{\E[\regret_T(\vecx_\star)]}{\weight_{1:T}} + \frac{D}{\weight_{1:T}}\sum_{t=1}^T \frac{(\sigma_G+D\sigma_H) t^{k-1/2}}{\sqrt{\lambda}} + \frac{\RDPdist(G+DH)\log(T) t^{k-1}}{\sqrt{\lambda}\rho}\\
    &\lesssim \frac{\E[\regret_T(\vecx_\ast)]}{T^{k+1}} + \frac{D}{\sqrt{\lambda}}\left(\frac{\sigma_G+D\sigma_H}{\sqrt{T}}+\frac{\RDPdist(G+DH)\log(T)}{\rho T}\right).
\end{align*}
For the second part of the theorem, 
\begin{align*}
    &\E[\|g_t+\gamma_t\|_*^2]
    \leq 3\E[\|\weight_t\nabla f(x_t)\|_*^2 + \|g_t-\weight_t\nabla f(x_t)\|_*^2 + \|\gamma_t\|_*^2] \\
    &\leq 3t^{2k}G^2 + \frac{12(k+1)^2(\sigma_G^2+D^2\sigma_H^2)t^{2k-1}}{\lambda}
    + \frac{24V(k+1)^2(G+DH)^2\log_2^2(2T)t^{2k-2}}{\lambda \rho^2} \\
    &= 3t^{2k}\left(G^2 + \frac{4(k+1)^2}{\lambda}\left(\frac{(\sigma_G+D\sigma_H)^2}{t}+\frac{2V(G+DH)^2\log_2^2(2T)}{\rho^2t^2}\right) \right).
\end{align*}
In the second line, the first term is bounded by Lipschitzness, the second term follows from \eqref{eq:DP-OTB:variance-norm}, and the third is bounded by the combination of \eqref{eq:tree-mechanism:noise} and \eqref{eq:DP-OTB:RDP-noise}, namely
\begin{equation}
    \E[\|\overline\vecr_t\|_\ast^2] \le \frac{8(k+1)^2(G+DH)^2\RDPdist^2\log_2(2T)t^{2k-2}}{\lambda\rho^2}.
    \label{eq:DP-OTB:main:total-tree-noise}
\end{equation}
We then conclude the proof by putting everything together.
\end{proof}